\subsection{\problemblock{*Egs} Data Block}\label{sec:block-egs}

The Engineering Graphics System (EGS) is available on the Silicon Graphics
workstations in the Aerospace Systems Development Center at Sandia.\manualref{4}
This software is designed to interactively plot sets of analysis data, such as
trajectory data. The data for EGS must be in a special format called the EGS database
format, so the \problemblock{*egs} data block is available to produce a file
containing trajectory variables in this format. It provides a simple interface to EGS,
so that high-quality plots of the trajectory data can be produced quickly.

The \problemblock{*egs} data block contains a filename and a list of variable names.
The file produced by this data block, an EGS database file, usually has the file type
\texttt{.dbf}. A database file is created for each \problemblock{*egs} data block,
and any number of \problemblock{*egs} data blocks can appear within a problem
file.

An example of an \problemblock{*egs} data block is

\begin{lstlisting}[style=taosmanual]
*egs traj.dbf
  time alt range vel mach long latgd gamgd
  psigd dynprs alpha nx ny nz ntotal
\end{lstlisting}

where the database filename is \texttt{traj.dbf} and 15 output variables have been
listed. Any number of output variables can be listed in the data block.
\Cref{fig:egs-database-example} contains a listing of the first part of the database
file created from this example.

\begin{center}
  \begin{minipage}{0.96\linewidth}
    \lstinputlisting[style=taosmanual,basicstyle=\ttfamily\scriptsize]{examples/chapter04/egs-database-example.txt}
    \captionof{figure}{Example of an EGS database file.}
    \label{fig:egs-database-example}
  \end{minipage}
\end{center}

Data is automatically written to the database file for each trajectory defined in the
problem, so no trajectory numbers are given. The units and output format of each
variable can be set with the \problemblock{*units/fmt} data block
(\cref{sec:block-units-fmt}).

TAOS overwrites the \problemblock{*egs} database files without warning each time
it runs. This automatically updates the files so they contain the latest trajectory
information for a problem. If this is not desirable, then the filename must be changed
between runs or the database files can be renamed before each new run.

Radar and relative-vehicle output variables, those starting with \texttt{rad} or
\texttt{rel}, must contain a subscript enclosed in square brackets. The subscript gives
the radar station number or the trajectory number for the relative-vehicle
calculations. For example, the variable \statevar{radrng[2]} is used to output the
radar range of each trajectory from radar station number 2, and the variable
\statevar{relvel[4]} is used to output the closing velocity of trajectory number 4
relative to the other trajectories.

If surveys have been defined (\cref{sec:block-survey}), the survey loops form
different data levels in the EGS database file so EGS can be used to plot the
trajectory data as a function of the survey variables. In
\cref{fig:egs-database-example} the variable \statevar{gama} is a survey-loop
variable, and a set of trajectories is computed for each value it takes on.

Summary variables (\cref{sec:block-summarize}) are not written to the standard
EGS database file. If summary variables have been defined, they can be written to
an EGS database file with

\begin{lstlisting}[style=taosmanual]
*egs summary file.dbf
\end{lstlisting}

The \texttt{summary} keyword and the filename are required; no output or summary
variable names are given. This form of the \problemblock{*egs} data block, with the
\texttt{summary} keyword and a filename, requires at least one survey loop and at
least one summary variable. All summary variables are written to the EGS database
file as a function of the survey variables.
